\chapter{Process Loading and Memory Layout}

\section{Introduction}

When you type \texttt{./myprogram} and press enter, what actually happens? How does a simple executable file on disk become a running process with memory, stack, and heap? This chapter explores the final piece of the compilation story: \textbf{process loading} and \textbf{memory layout}.

We've seen how source code becomes an executable file. Now we'll see how that executable file is brought to life by the operating system's loader. We'll examine:
\begin{itemize}
    \item \textbf{The loader}: How the OS reads and maps executable files
    \item \textbf{Virtual memory}: How each process gets its own address space
    \item \textbf{Memory segments}: text, data, bss, heap, stack, and more
    \item \textbf{Dynamic linking}: How shared libraries are loaded at runtime
    \item \textbf{Memory mapping}: How \texttt{mmap} enables efficient file I/O and memory allocation
\end{itemize}

Understanding process loading is crucial for:
\begin{itemize}
    \item \textbf{Debugging}: Understanding segmentation faults, memory errors
    \item \textbf{Performance}: Optimizing memory usage and allocation strategies
    \item \textbf{Security}: Understanding ASLR, stack protection, and memory isolation
    \item \textbf{Systems programming}: Writing low-level code that interacts with the OS
\end{itemize}

\section{From Executable to Process}

\subsection{What is a Process?}

A \textbf{process} is an instance of a program in execution. Unlike a program (static file on disk), a process is dynamic:

\begin{verbatim}
Program (File on Disk)          Process (Running Instance)
├── Code sections              ├── Memory regions
├── Data sections              ├── Stack
├── Metadata (headers, etc.)   ├── Heap
└── ...                        ├── Program counter
                               ├── Open files
                               └── ...
\end{verbatim}

\textbf{Key differences}:
\begin{itemize}
    \item \textbf{Program}: Static, passive, stored on disk
    \item \textbf{Process}: Dynamic, active, lives in memory
\end{itemize}

\subsection{The Loading Process}

When you run a program, the OS performs these steps:

\begin{verbatim}
1. User executes: ./myprogram
         │
         ▼
2. Shell calls fork() + execve()
         │
         ▼
3. Kernel's execve() handler:
   a. Read ELF header from file
   b. Verify it's a valid executable
   c. Create new memory map for process
   d. Read program headers (PHDRs)
   e. Map segments into memory (mmap)
   f. Set up stack and arguments
   g. Set up dynamic linker (if needed)
   h. Jump to entry point (_start)
         │
         ▼
4. _start (in libc/crt0):
   a. Initialize libc
   b. Run constructors (.init_array)
   c. Call main()
   │
   ▼
5. Your program runs!
   │
   ▼
6. After main() returns:
   a. Run destructors (.fini_array)
   b. Call exit()
   │
   ▼
7. Kernel cleans up process
\end{verbatim}

\section{Virtual Memory}

\subsection{What is Virtual Memory?}

\textbf{Virtual memory} gives each process the illusion of having its own private, contiguous address space. In reality, physical memory is fragmented and shared among all processes.

\textbf{With virtual memory}:
\begin{itemize}
    \item Each process has its own address space (0 to $2^{64}-1$ on 64-bit)
    \item The Memory Management Unit (MMU) translates virtual $\rightarrow$ physical addresses
    \item Processes are isolated from each other
    \item Pages can be moved, swapped, or shared transparently
\end{itemize}

\textbf{Translation process}:
\begin{verbatim}
Virtual Address (seen by program)
        │
        ▼
MMU (Memory Management Unit)
        │
        ├──► Page Table Lookup
        │
        ▼
Physical Address (actual RAM location)
\end{verbatim}

\subsection{Benefits of Virtual Memory}

\begin{enumerate}
    \item \textbf{Isolation}: Processes can't access each other's memory
    \item \textbf{Protection}: Pages can be marked read-only, execute-only, etc.
    \item \textbf{Efficiency}: Only loaded portions of a program need to be in RAM
    \item \textbf{Simplicity}: Each program can use the same address range
    \item \textbf{Sharing}: Multiple processes can share the same physical page (for libraries)
\end{enumerate}

\section{Memory Segments}

\subsection{Typical Process Memory Layout}

A typical process memory layout on x86-64 Linux:

\begin{verbatim}
High Addresses (0xFFFFFFFFFFFFFFFF)
     ┌─────────────────────────────┐
     │      Kernel Space           │  (Not accessible to user)
     ├─────────────────────────────┤ 0x7FFFFFFFF000
     │                             │
     │      Stack                  │  Grows downward
     │      ┌─────────────┐        │
     │      │  stack var  │        │
     │      ├─────────────┤        │
     │      │     ...     │        │
     │      └─────────────┘        │
     │                             │
     ├─────────────────────────────┤ ← RSP (stack pointer)
     │                             │
     │          ↓                  │  Free memory
     │                             │
     │          ↑                  │
     │                             │
     ├─────────────────────────────┤
     │      Memory Mapping Segment │  Libraries, mmap regions
     │      ┌─────────────┐        │
     │      │   shared lib│        │
     │      ├─────────────┤        │
     │      │     ...     │        │
     │      └─────────────┘        │
     ├─────────────────────────────┤
     │      Heap                   │  Grows upward (via brk/mmap)
     │      ┌─────────────┐        │
     │      │  malloc'd   │        │
     │      ├─────────────┤        │
     │      │     ...     │        │
     │      └─────────────┘        │
     │                             │
     ├─────────────────────────────┤ ← Program break (brk)
     │      .bss                   │  Uninitialized data
     │      (0-filled at load)     │
     ├─────────────────────────────┤
     │      .data                  │  Initialized data
     │      (global vars, etc.)    │
     ├─────────────────────────────┤
     │      .rodata                │  Read-only data (strings, etc.)
     ├─────────────────────────────┤
     │      .text                  │  Code (executable)
     │      (machine instructions) │
     └─────────────────────────────┘
Low Addresses (0x0000000000000000)
\end{verbatim}

\subsection{The .text Segment}

Contains executable code (machine instructions).

\textbf{Properties}:
\begin{itemize}
    \item \textbf{Read-only}: Code cannot modify itself
    \item \textbf{Executable}: CPU can fetch and execute instructions
    \item \textbf{Shared}: Multiple instances of the same program can share this segment
\end{itemize}

\subsection{The .data Segment}

Contains \textbf{initialized} global and static variables.

\textbf{Properties}:
\begin{itemize}
    \item \textbf{Read-write}: Variables can be modified
    \item \textbf{Initialized}: Contains initial values from executable
    \item \textbf{Fixed size}: Size known at compile time
\end{itemize}

\subsection{The .bss Segment}

Contains \textbf{uninitialized} global and static variables.

\textbf{Properties}:
\begin{itemize}
    \item \textbf{Read-write}: Variables can be modified
    \item \textbf{Zero-initialized}: Set to 0 at process start
    \item \textbf{Not stored in executable}: Only size is recorded, not data
\end{itemize}

\textbf{Why separate .bss from .data?}

Efficiency! .bss doesn't take up space in the executable file:

\begin{verbatim}
.data in file:
[initial values for 100 bytes] → 100 bytes in executable

.bss in file:
[size: 1000, zero at load] → 8 bytes in executable (just the size!)
\end{verbatim}

The loader allocates .bss memory and zeros it at runtime.

\subsection{The Heap}

Dynamic memory allocation region managed by \texttt{malloc}/\texttt{free}.

\textbf{Growth direction}: Upward (toward higher addresses)

\textbf{Management}:
\begin{itemize}
    \item libc's heap allocator (ptmalloc, jemalloc, etc.) manages this region
    \item Requests more memory from kernel via \texttt{brk()} or \texttt{mmap()} when needed
    \item Carves large regions into small allocations for your program
\end{itemize}

\subsection{The Stack}

Function call region for local variables, return addresses, saved registers.

\textbf{Growth direction}: Downward (toward lower addresses)

\textbf{Stack overflow}: Exceeding stack size causes segmentation fault

\section{The Loader in Detail}

\subsection{Program Headers}

ELF executables contain \textbf{program headers} that tell the loader how to load the file.

\textbf{Key program headers}:
\begin{itemize}
    \item \textbf{PHDR}: Program header table itself
    \item \textbf{INTERP}: Specifies dynamic linker (ld-linux.so)
    \item \textbf{LOAD}: Loadable segment (code or data)
    \item \textbf{DYNAMIC}: Dynamic linking information
    \item \textbf{GNU\_STACK}: Stack permissions (executable? non-executable?)
\end{itemize}

\subsection{The Dynamic Linker}

For dynamically linked executables, the \textbf{dynamic linker} (\texttt{ld-linux.so}) performs final setup:

\textbf{What it does}:
\begin{enumerate}
    \item Load required shared libraries (libc.so, etc.)
    \item Resolve symbol references (PLT/GOT, as seen in Chapter 7)
    \item Run relocations
    \item Call shared library initializers
    \item Transfer control to the program
\end{enumerate}

\subsection{ASLR: Address Space Layout Randomization}

Modern systems use \textbf{ASLR} to randomize memory locations for security:

\textbf{Without ASLR} (old systems):
\begin{verbatim}
Stack always at 0x7ffffff00000
Heap always starts at 0x02000000
Libraries always loaded at same addresses
→ Predictable = easier to exploit!
\end{verbatim}

\textbf{With ASLR} (modern systems):
\begin{verbatim}
Stack at random location each run
Heap starts at random location
Libraries loaded at random addresses
→ Unpredictable = harder to exploit!
\end{verbatim}

\section{Key Takeaways}

\begin{enumerate}
    \item \textbf{Loading transforms executables into processes}: The OS loader reads the ELF file, maps segments into memory, and transfers control to the entry point.
    \item \textbf{Virtual memory isolates processes}: Each process has its own address space, translated to physical memory by the MMU via page tables.
    \item \textbf{Memory is organized into segments}: .text (code), .rodata (read-only data), .data (initialized data), .bss (uninitialized data), heap, and stack.
    \item \textbf{The stack grows downward, heap grows upward}: Local variables and function frames are on the stack; dynamic allocations are on the heap.
    \item \textbf{The loader uses program headers}: ELF program headers specify which segments to load, their permissions, and their virtual addresses.
    \item \textbf{Dynamic linking happens at load time}: The dynamic linker loads shared libraries, resolves symbols, and performs relocations before your code runs.
    \item \textbf{ASLR randomizes memory layout}: For security, ASLR randomizes stack, heap, and library locations to make exploitation harder.
\end{enumerate}

\section{Looking Ahead}

This completes our journey from C source code to running process! We've seen:
\begin{itemize}
    \item \textbf{Chapter 1-2}: Compilation pipeline and preprocessing
    \item \textbf{Chapter 3-4}: Compilation internals and object files
    \item \textbf{Chapter 5-6}: ELF format and linking
    \item \textbf{Chapter 7-8}: Static/dynamic libraries and system calls
    \item \textbf{Chapter 9}: ABI and cross-platform compilation
    \item \textbf{Chapter 10}: Process loading and memory layout
\end{itemize}

You now understand the complete path from source code to executing process, with deep knowledge of each step along the way.
